\scanpage{59}
\mspage{II. 7}
% lines 1-6
$A$ est décrit par une représentation
\[
\Gamma\mathrel{\text{---}}\operatorname{Aut}(B)
\]
où $\Gamma$ est le groupe de Galois de $X'/X$. L'hypothèse \unsure{indique|implique} que $\Gamma$ opère trivialement sur ${}_n B(k)$, et on sait \unsure{q|que} (comme $n\geq3$) que cela \unsure{implique|entraîne} que l'action de $\Gamma$ sur $B$ est triviale, d'où la conclusion.

% lines 7-8
\uline{Remarque} 1.8. \deleted{Le théorème \unsure{précédent}} Par descente, on voit que 1.7 reste valable tel quel sans supposer $k$ \deleted{alg.} \inserted{sép.} clos.

% lines 9-12
L'\uline{énoncé} 1.5 est de nature purement géométrique, (bien que sa démonstration soit de nature arithmétique). Voici une variante arithmétique de 1.5 :

% lines 13-23
\uline{Théorème} 1.9. Soit $K$ un corps de type fini \deleted{\unsure{$A$}} $\overline K$ une clôture séparable de $K$, \inserted{$\pi=\operatorname{Gal}(\overline K/K)$,} $A_K$ un schéma abélien sur $K$. Supposons l'action de $\pi$ sur $T_\ell(A(\overline K))$ \deleted{abélienne} \inserted{essentiellement résoluble}. Alors, il existe une extension finie \deleted{\unsure{$B$}} \inserted{$k'\subset\overline K$} du corps premier, \deleted{et un} \inserted{et} un schéma abélien $B$ sur $k'$, et une isogénie radicielle \deleted{$B_{K'}\simeq A$}
\[
B\otimes_{k'}K'\simeq A\otimes_K K',
\]
et on a posé \deleted{\unsure{$\ell$}} $K'=Kk'$. Lorsque $\operatorname{car}K=0$, \deleted{\unsure{donc}} \inserted{l'action de $\pi$ \deleted{\unsure{sur}} sur $T_\ell(A(\overline K))$ est \uline{ess\textsuperscript{t} abélienne} $(*)$ ou tout au moins \unsure{ess\textsuperscript{t} nilpotente}} $k'$ est \deleted{de car nulle} un corps de nbs, on peut supposer que $B$ \deleted{\unsure{est} \unsure{$R'$} \unsure{fini}} provient d'un schéma abélien sur \inserted{$R$ \deleted{\unsure{$R'$}}} l'anneau des entiers de $k'$.

% footnote
$(*)$ Hypothèse sans doute \deleted{\unsure{conséquence}} \unsure{autom. satisfaite} \unsure{Donc à|Donc on} \unsure{démontrer que} l'action est tjs $\frac12$ simple.
% archive sheet number: 58
